É possível ver este movimento num trecho de Krenak, pp. 56-57: \enquote{ainda segundo Viveiros de Castro, o Brasil se especializou na produção de pobres. Nossa tecnologia para produzir pobreza é mais ou menos assim: a gente pega quem pesca e colhe frutos nativos, tira do seu território e joga nas periferias da cidade, onde nunca mais vai poder pegar um peixe para comer, porque o rio que passa no bairo está podre. Se você tira um Yanomami da floresta, onde ele tem água, alimento e autonomia, e bota em Boa Vista, isso é produção de pobreza. Se expulsa o pessoal da Volta Grande do Xingu para fazer uma hidrelétrica, mandando para um beiradão de Altamira, você está convertendo-o em pobre. O capitalismo precisa de uma plataforma -- que é urbana.} Mais do que a urbanidade, diria que a plataforma capitalista é a escassez material que faz do dinheiro um articulador de tudo necessário para praticamente todas as dimensões da vida.